\SourceLeaf{285}%
\begin{center}\begin{tabular}{r|ccc}
Anfangspunkt&&&\\
&$-a$&$-b$&$+c$\\
8&&&\\
&$-a$&&\\
9&&&\\
&$+a$&&\\
12&&&\\
&$-a$&&\\
14&&&\\
&$+a$&&\\
16&&&\\
&$+a$&$-b$&\\
17&&&\\
&$-a$&$-b$&\\
18&&&\\
&$+a$&$-b$&\\
19&&&\\
&$+a$&$-b$&$+c$\\
20&&&\\
&$+a$&$+b$&$+c$\\
21&&&\\
&$-a$&$+b$&$+c$\\
22&&&\\
&$-a$&$-b$&$+c$\\
Anfangspunkt&&&\\
\end{tabular}\end{center}
Man sieht aus der vorigen Analyse, dass das letzte Zeichen, welches in
diesem Tableau jeder Figur entspricht, nothwendig dasselbe sein muss,
wie das erste, oder, dass die Zeichen $+-$, die bei den einzelnen
Intersectionen vorkommen, wenn man diejenigen abrechnet, die bei den
einzelnen Bildungspunkten der Figuren vorkommen, einander destruiren;
die Summe \textls{aller} ist folglich dieselbe, wie die der letztern allein, und
zugleich einerlei mit der Summe der Zeichen in dem Character des
Anfangspunktes gegen die einzelnen Figuren $a$, $b$, $c$, d.\,i. alle,
wenn man die erst zuletzt gebildete $d$ nicht mitzählt.
\par\Needspace{12\baselineskip}
Der Character auf der einen Seite der Figur $d$ gegen die Einzelnen
Figuren werde durch $\alpha a+\beta b+\gamma c+\delta d$ aus[gedrückt]
und [man] setze $\alpha+\beta+\gamma+\delta=\sigma$, der absolute
Character jenes Platzes hingegen sei $=\chi$; dann ist die ganze
Seitenbiegung
\[
= -\sigma\,.\,360^0 - \frac{F-\chi S}{S}\,.\,720^0,
\]
